% Generated draft: requires mathematical and source review.
\documentclass[11pt]{article}
\usepackage[margin=25mm]{geometry}
\usepackage{fontspec}
\setmainfont{Latin Modern Roman}
\usepackage{amsmath,amssymb,mathtools,unicode-math}
\setmathfont{Latin Modern Math}
\usepackage{enumitem,needspace,xurl,hyperref}
\setlength{\parindent}{0pt}
\setlength{\parskip}{5pt}
\setlength{\emergencystretch}{3em}
\newcommand{\N}{\mathbb{N}}
\newcommand{\Z}{\mathbb{Z}}
\newcommand{\Q}{\mathbb{Q}}
\newcommand{\R}{\mathbb{R}}
\newcommand{\C}{\mathbb{C}}
\newcommand{\F}{\mathbb{F}}
\newcommand{\A}{\mathbb{A}}
\newcommand{\PP}{\mathbb{P}}
\newcommand{\E}{\mathbb{E}}
\newcommand{\eps}{\varepsilon}
\newcommand{\dd}{\,\mathrm{d}}
\begin{document}
\section*{UnsolvedMath ID: 30006512}
\textbf{Source problem code:} catalogue entry
\subsection{Definitions and mathematical statement}
The supplied record asks whether the finite-domain constraint-satisfaction dichotomy extends to reducts of finitely bounded homogeneous structures. It does not explicitly specify the signature restrictions, the meaning of reduct, or the complexity reductions. The following is a precise candidate formulation under conventional relational interpretations; these choices require source confirmation.

A relational signature is a collection of relation symbols, each with a specified finite arity. A relational structure is homogeneous if every isomorphism between two finite induced substructures extends to an automorphism of the whole structure. For a relational structure \(\mathbb A\), its age \(\operatorname{Age}(\mathbb A)\) is the class, up to isomorphism, of its finite induced substructures. Here, finitely bounded means that there is a finite family \(\mathcal F\) of finite structures in the signature of \(\mathbb A\) such that, for every finite structure \(\mathbb C\) in that signature,
\[
\mathbb C\in\operatorname{Age}(\mathbb A)
\quad\Longleftrightarrow\quad
\text{no }\mathbb F\in\mathcal F\text{ embeds into }\mathbb C
\text{ as an induced substructure}.
\]
A parameter-free first-order reduct \(\mathbb B\) of \(\mathbb A\) is a relational structure on the same domain whose basic relations are definable in \(\mathbb A\) by first-order formulas without parameters; its signature need not be a subsignature of that of \(\mathbb A\).

For a fixed relational structure \(\mathbb B\) of finite signature, \(\operatorname{CSP}(\mathbb B)\) is the following decision problem: given a finite structure \(\mathbb X\) in the same signature, determine whether there exists a homomorphism \(h:\mathbb X\to\mathbb B\). Thus, for each relation symbol \(R\) of arity \(r\) and every \((x_1,\ldots,x_r)\in R^{\mathbb X}\), one requires
\[
(h(x_1),\ldots,h(x_r))\in R^{\mathbb B}.
\]
The template \(\mathbb B\) is fixed and only \(\mathbb X\) is input.

\textbf{Candidate dichotomy question.} For every nonempty countable homogeneous relational structure \(\mathbb A\) in a finite signature whose age is finitely bounded, and every parameter-free first-order reduct \(\mathbb B\) of \(\mathbb A\) in a finite relational signature, is it true that
\[
\operatorname{CSP}(\mathbb B)\in\mathrm P
\quad\text{or}\quad
\operatorname{CSP}(\mathbb B)\text{ is }\mathrm{NP}\text{-complete}
\]
under polynomial-time many-one reductions? Here countable permits both finite and countably infinite domains. The alternatives are not asserted to be mutually exclusive without an additional complexity assumption.

The countability convention, finite signatures, parameter-free first-order interpretation of reduct, induced-embedding interpretation of finite bounds, and reduction convention are explicit interpretive choices, not hypotheses recovered from the supplied excerpt. The excerpt merely mentions an analogous conjecture for homogeneous structures; it does not supply a precise algebraic characterization or the surrounding conjecture. Consequently, the candidate above requires human mathematical and source review before it can be identified with the exact intended conjecture.
\subsection{Short English statement}
Under the stated conventional interpretation, the question asks whether every finite-signature constraint-satisfaction problem obtained from a countable finitely bounded homogeneous relational structure by parameter-free first-order definitions is either polynomial-time decidable or \(\mathrm{NP}\)-complete. The supplied excerpt leaves several hypotheses and the exact intended dichotomy unspecified. The record labels the problem open; this is supplied curation, not an independently verified status.
\subsection{Sources}
\begin{itemize}
\item Supplied local UnsolvedMath registry record, permanent ID \(30006512\), title \emph{Constraint-Satisfaction Dichotomy for Homogeneous Structures}, exact source problem code \texttt{catalogue entry}; collection: Oberwolfach Reports Open Problems. The supplied proposed year is 2025; no individual proposer is named.
\item Supplied source citation: \emph{Homogeneous Structures: Model Theory meets Universal Algebra} (2026), Oberwolfach Reports, pp. 3063--3124, DOI \texttt{10.4171/owr/2025/58}. Supplied source URL: \url{https://doi.org/10.4171/owr/2025/58}. This report is the supplied authoritative reference; its contents have not been independently consulted for this draft.
\item Supplied literature reference: \emph{Oberwolfach Report 2025/58}, \url{https://doi.org/10.4171/OWR/2025/58}. The exact supplied literature-check date is \texttt{2026-08-22}. The assessment says that the finite-domain constraint-satisfaction dichotomy is known and that an analogous dichotomy for homogeneous structures remains open. This is supplied dated curation derived from the OWR record, not an independently verified literature assessment; time-sensitive status and cited claims require independent verification.
\item The registry supplies the statement-status label \texttt{corrected\_verified} and reports that the official report was checked, the question isolated from surrounding exposition, and notation repaired. These are source-provided verification claims, not verification performed for this draft. The supplied original excerpt states only that an analogous conjecture for homogeneous structures remains open despite structural insights.
\item The original excerpt gives Mathematics Subject Classification (2020): 03Cxx, 08Axx, and states that, unless otherwise noted, report content is licensed under CC BY SA 4.0. The supplied report-title year 2026 differs from the 2025 report identifier and proposed year; the chronology is retained without reconciliation.
\end{itemize}
\end{document}
